\documentclass[12pt]{scrreprt}

\usepackage[ngerman]{babel}
\usepackage{textcomp}
\usepackage[T1]{fontenc}
\usepackage[utf8]{inputenc}

\usepackage{color,fancyvrb}
\usepackage{lmodern}

\usepackage{graphicx}
\usepackage{amssymb,amsmath}
%\graphicspath{{../latex/}}

%% Quellcode auszeichnung %%
\usepackage{minted}
\makeatletter
\minted@define@extra{label}
\makeatother
\definecolor{bg}{rgb}{0.95,0.95,0.95}
\usemintedstyle{manni}
\newminted{c}{fontfamily=helvetica, fontsize=\scriptsize, tabsize=8,
samepage, frame=single, bgcolor=bg,label=Quellcode,
numbersep=5pt, xrightmargin=0mm, xleftmargin=4mm, linenos}

\usepackage{url}
\usepackage{wrapfig}
\usepackage[font=small,labelfont=bf,textfont=sf]{caption}
\usepackage{placeins}
\usepackage{multirow}
\setlength{\parskip}{3ex plus 2ex minus 2ex}
\widowpenalty=300
\clubpenalty=300

\newcommand{\comment}[1]{}
\setcounter{secnumdepth}{6}

% header and footer
\usepackage{fancyhdr}
\setlength{\headheight}{18pt}
\renewcommand{\headrulewidth}{0.3pt}
\pagestyle{fancyplain}
\renewcommand{\chaptermark}[1]{\markboth{#1}{}}
\fancyhf{}
\cfoot{\thepage}

% make internal and external links
\usepackage{hyperref} % load this at the end of all packages, but
%before the options, not true here at the moment tough

\hypersetup{
%    bookmarks=true,         % show bookmarks bar? - is already true
    unicode=false,          % non-Latin characters in Acrobat’s bookmarks
    pdftoolbar=true,        % show Acrobat’s toolbar?
    pdfmenubar=true,        % show Acrobat’s menu?
    pdffitwindow=false,     % window fit to page when opened
    pdfstartview={FitH},    % fits the width of the page to the window
    pdftitle={CG 1 - Exercise},    % title
    pdfauthor={CG 1 - Gruppe 6},     % author
    pdfsubject={},   % subject of the document
    pdfcreator={CG 1 - Gruppe 6},   % creator of the document
    pdfproducer={Producer}, % producer of the document
    pdfkeywords={CG 1} {Computer Graphics I} {TU Berlin}, % list of keywords
    pdfnewwindow=true,      % links in new window
    colorlinks=false,       % false: boxed links; true: colored links
    linkcolor=red,          % color of internal links
    citecolor=green,        % color of links to bibliography
    filecolor=magenta,      % color of file links
    urlcolor=cyan           % color of external links
}

\begin{document}

\begin{titlepage}
  \begin{center}
    \vspace*{1cm}
    \textbf
    {\emph
      {\Large Computer Graphics I}}\\[1cm]
    \textbf
    {\Large Technische Universität Berlin}\\[.5cm]
    \textbf
    {\Large WS 2012/2013}\\[2cm]
    %\includegraphics[bb=0 0 60 40]{./TUB_Logo.png}\\[0.9cm]
    \vfill
    {\Large Gruppe 11}\\[1cm]
    Christopher Sierigk\\
    Bernd Loeber\\
    Silvio Hädrich\\
    \vspace*{2cm}
    \today
  \end{center}
  \vspace*{2cm}
\end{titlepage}

\pagenumbering{arabic}

\chapter*{Aufgabe 1 - Praktischer Teil}

siehe Quellcode

\chapter*{Aufgabe 2 - Theoriefragen}
\section*{1.}
Der Grund für das Clipping am kanonischen Sichtvolumen in OpenGL ist Performance. Durch das Clipping werden alle Teile der Geometrie abgeschnitten, die außerhalb des sichtbaren Bereichs liegen. Dies hat den Grund beim Rendern Rechenarbeit zu sparen, da es wenig Sinn macht Objekte bzw. Teile von Objekten zu rendern, die eindeutig nicht gesehen werden können.
\section*{2.}
Da ein 3dimensionales kanonisches Sichtvolumen 6 Seiten hat, besteht der Outcode aus 6 Ziffern die entweder 1 oder 0 sein können. Die Seiten außerhalb des Volumens kriegen jeweils einen eindeutigen Code zugeordnet:

\begin{align*}
	000001& &top\\
	000010& &bottom\\
	000100& &left\\
	001000& &right\\
	010000& &front\\
	100000& &back
\end{align*}

Wenn sich 2 Seiten an einer Kante schneiden werden die jeweiligen Codes verodert:

\begin{align*}
	000101& &top-left\\
	001001& &top-right\\
	010001& &top-front\\
	100001& &top-back\\
	000110& &bottom-left\\
	001010& &bottom-right\\
	010010& &bottom-front\\
	100010& &bottom-back\\
	010100& &front-left\\
	011000& &front-right\\
	100100& &back-left\\
	101000& &back-right
\end{align*}

Treffen sich die Kanten an einem Eckpunkt entstehen wiederum verschiedene Outcodes:

\begin{align*}
	0100101& &top-left-front\\
	0110001& &top-right-front\\
	1001001& &top-left-back\\
	1010001& &top-right-back\\
	0100110& &bottom-left-front\\
	0110010& &bottom-right-front\\
	1001010& &bottom-left-back\\
	1010010& &bottom-right-back
\end{align*}

Als Primitive eignen sich im 3D-Fall die Ebenen.

\section*{3.}
\section*{4.}
Die Komplexität bezieht sich dabei auf das "Primitiv" der Dreiecke.
\section*{5.}
\section*{6.}
\subsection*{(a)}
Wir nehmen einen beliebigen Vektor $\vec{v}$ aus der Ebene und die Normale $\vec{n}$ der Ebene. Dabei lässt sich die Ebene wie folgt beschreiben:

\[
	\vec{n} \cdot \vec{v} = n_x \cdot v_x + n_y \cdot v_y + n_z \cdot v_z = 0
\] 
\subsection*{(b)}
\subsection*{(c)}
Das Transformationsgesetz bei einer Normalen gilt aus dem Grund nicht, da eine Normale nicht durch eine Punkt zu Punkt Beziehung beschrieben werden kann. Das bedeutet, es lässt sich keine Normale nur aus 2 gegebenen Punkten bestimmen und somit verhält sich eine Normale auch nicht wie ein zu transformierender Punkt.
\subsection*{(d)}
\end{document}
